\section{Classification results and scope}\label{sec:results}
A fusion ring is a finite-rank based ring with nonnegative integral
structure constants, a distinguished unit, and duality satisfying
Frobenius reciprocity. We work with full ordered multiplication tensors;
in particular, commutativity is not part of the definition.
The multiplicity is the maximum of all structure constants, including
the unit coefficients. A categorification means a fusion category over
$\CC$, without an assumed pivotal, spherical, or unitary structure.
We classify existence for a ring, not equivalence classes of categories
with that ring.

Vercleyen--Slingerland constructed a substantial low-rank database
\cite{VS}. Its rank-six portion is complete through multiplicity four.
The multiplicity-one categorification problem was settled by
Liu--Palcoux--Ren \cite{LPR}. The rank-four and rank-five projects
\cite{DP4,DP5} develop arithmetic, sphericalization, and central-induction
certificates. At rank six, however, neither the ring nor a signed quotient
of its sphericalization can simply be assumed commutative. This is the
principal structural change addressed here.


All positively identified rings and their models appear in
Table~\ref{tab:models}; its structure column is part of the classification
statement, not merely a description of a chosen presentation. In
particular, Schopieray's category \cite[Example~1]{SchNonPU} realizes
$\RR{2}{124}$ spherically, whereas \cite[Theorem~1]{SchNonPU} excludes
\emph{every} pseudo-unitary categorification of that same ring.
Thus general or spherical existence cannot be identified with unitary
existence, even in rank six and multiplicity two. The electronic ledger
records general, pivotal, spherical, pseudo-unitary and unitary existence
separately for every tensor. The source identifications are detailed in Section~\ref{sec:literature-models}.
The paper is organized by mathematical operation rather than by individual
exceptions. Section~\ref{sec:enumeration} fixes the source census, its scope, and
the tensor conventions used by every automated test. Section~\ref{sec:reductions} reduces arbitrary
categorifications to the appropriate arithmetic and spherical cases.
Section~\ref{sec:unified-induction} presents one central-induction framework,
with precisely scoped refinements rather than separate ad hoc tests.
Section~\ref{sec:newalgebra} supplies module and Morita constraints.
Section~\ref{sec:positive} treats models by construction, including
central-boson quotients, equivariantizations, and relative centers.
Section~\ref{sec:certification} specifies exhaustive certificate checking
and the remaining frontier.

In particular, the positive catalogue includes the extended Haagerup
category $\mathcal{EH}_1$ \cite{BMPS,BPEH} and Izumi's $C_5+5$ near-group
category \cite{Izumi}, not just the constructions derived from them.
The mixed fermionic quotient of Proposition~\ref{prop:mixed-fermion}, the
sign-swap equivariantization of Theorem~\ref{thm:sign-swap}, and the
relative-center construction of Corollary~\ref{cor:tt3-relative} give
three additional full-tensor identifications. None is inferred merely
from a dimension match or from survival of the necessary criteria.

\begin{theorem}[Source census {\cite{CensusRepo}}]\label{thm:census}
The published \emph{Fusion Ring Census}, at the revision fixed in
Section~\ref{sec:enumeration}, contains exactly $16612$
based-isomorphism classes of rank-six fusion rings of multiplicity at
most eight. The numbers at exact multiplicities $1,\ldots,8$ are
\[
39,\quad154,\quad384,\quad872,\quad1582,\quad2768,\quad3814,\quad6999.
\]
All three possible duality types and all noncommutative cases are
included. No categorifiability filter is applied. Completeness of this
counting input is the result of the cited repository; the present paper
addresses categorification of its tensors.
\end{theorem}

\begin{theorem}[Classification through multiplicity two and certified extension]\label{thm:classification}
In the census of Theorem~\ref{thm:census}, the accompanying exact
certificates exclude $\NegativeCount$ rings from general complex
categorification. There are $\PositiveCount$ positively identified rings:
$\UnitaryCount$ have unitary categorifications, whereas $\RR{2}{124}$ has
a spherical categorification but no pseudo-unitary categorification.
These positive sets and the general exclusions are pairwise disjoint.
The remaining $\UnresolvedCount$ rings are not decided for general existence by the present criteria and bounded searches. Structure-specific pseudo-unitary exclusions are recorded separately and do not alter that general status. Their distribution is given in
Table~\ref{tab:classification}.

Among the $193$ rings of multiplicity at most two, precisely $43$
admit general, pivotal, or spherical categorifications. Exactly $42$ of
these admit pseudo-unitary or unitary categorifications; the sole
exception is $\RR{2}{124}$. At exact multiplicity two the $154$ rings
split into $132$ general exclusions, $21$ unitary positives and that
one spherical non-pseudo-unitarily-categorifiable ring. There are no
residuals at multiplicity two. The higher-multiplicity results in
Table~\ref{tab:classification} are partial; their residuals are listed
in the complete electronic ledger.
\end{theorem}

\begin{table}[ht]\centering\small
\begin{tabular}{rrrrrr}\toprule
Exact multiplicity & Rings & Excluded & Unitary & S, not PU & Undecided\\\midrule
1 & 39 & 18 & 21 & 0 & 0\\
2 & 154 & 132 & 21 & 1 & 0\\
3 & 384 & 380 & 3 & 0 & 1\\
4 & 872 & 855 & 2 & 0 & 15\\
5 & 1582 & 1577 & 1 & 0 & 4\\
6 & 2768 & 2762 & 3 & 0 & 3\\
7 & 3814 & 3807 & 2 & 0 & 5\\
8 & 6999 & 6980 & 0 & 0 & 19\\
\midrule
Total & 16612 & 16511 & 53 & 1 & 47\\
\bottomrule\end{tabular}
\caption{Complete counting census; categorification is complete through multiplicity two and partial above. The column S, not PU denotes a spherical realization with no pseudo-unitary categorification. All excluded rings are ruled out generally.}\label{tab:classification}
\end{table}

\paragraph{The low-multiplicity residuals.}
Every ring at exact multiplicity 2 is decided.
At exact multiplicity 3 the unresolved records are
\[\begin{gathered}
\RR{3}{215}.
\end{gathered}\]


The phrase ``undecided'' is deliberately local to this computation. A
search timeout is not an obstruction, and failure to produce a separating
functional is not a construction. Nor do we claim that every residual
is open throughout the literature. Section~\ref{sec:frontier} retains exact feasible points for the now positively identified cyclic-four ring as regression controls. Such points were not used to infer existence; the conformal-embedding construction supplies it.

For reproducible identifiers, retain the unit as the zeroth basis element,
use the exported representatives in the generator's fixed coordinate
conventions, and sort their full tensors in row-major lexicographic order
at each exact multiplicity. This order is a convention on the exported
representatives, not a claim that each is the lexicographic minimum under
all $120$ relabellings. Distinctness up to based isomorphism is supplied
by the cited census.
The established representatives and identifiers through multiplicity seven are preserved; the multiplicity-eight stratum is sorted by the same row-major convention. The $j$th tensor at multiplicity $m$ is denoted
$\RR{m}{j}$. The complete tensor, rather than this identifier, is the input
to every obstruction. A separate crosswalk records matches to the
ancillary database of \cite{VS}; its indices play no role in a proof.

All reductions are applied to full tensors. The cheap arithmetic filters
are retained because they avoid unnecessary induction searches; the
stronger blocks are used only with their stated hypotheses. The
rank-four and rank-five results are cited rather than reproved. This is
a unified implementation of sufficient exclusion certificates, not a
claim that the criterion set is logically minimal.

\paragraph{Progress at the first unresolved multiplicity.}
The earliest general residual remains $\RR{3}{215}$. We do not skip it
in favor of larger-multiplicity examples. Section~\ref{sec:joint-low-product-gram}
proves a joint low-product Gram condition, combining fourth-order
modular moments with integral product restrictions. Its exact automated
test shows that the formal rank-$22$ witness survives; this is a precise
limit of the strengthened necessary system, not a new general exclusion.
Section~\ref{sec:premodular-frontier-bound} derives from
Dong--Schopieray a rank-six bound $\FPdim\leq168$ for braided spherical
categorifications. Since the tensor has
$\FPdim=136+32\sqrt{17}>168$, it is excluded in that scope only.
A full-census tensor comparison with the rank-six ancillary data of
Ng--Rowell--Wen gives no new model. Each numerical conclusion is produced
by the supplied scripts from tensors and exact source data, without
manual case arithmetic.

\paragraph{What the exclusion table measures.}
The sequential table in Section~\ref{sec:verification} applies the
available criteria in an explicit order and reports new removals and
remaining tensors at each exact multiplicity. Simple structural and
arithmetic tests precede finite induction and signed searches. This is
a reproducible compromise between cost and yield, not a proof of a
universally optimal order. A theorem with a stronger conclusion can be
more expensive to apply, and a successful earlier test makes later tests
unnecessary for that tensor.
Dimension-order and character-induction stages first make additional
top-level removals at multiplicity eight, one and four respectively.
This describes the executed sieve through eight; it does not claim
that the same criteria decide larger multiplicities.

The earliest unresolved multiplicity is the first stratum in which the
current automated procedure does not decide every ring. It marks where
further work should begin; that work may be a new exclusion criterion,
a sharper application of an existing criterion, or a positive model.
A zero residual count in a later stratum means that the present methods
suffice for that stratum. Even if an entire terminal interval through
multiplicity eight has zero residuals, this would not imply that no new
criterion is needed at multiplicity nine or higher. Conversely, a
bounded search that stops without a certificate does not show that a
new theorem is logically necessary.

The same reporting scheme applies at lower ranks: one fixes a complete
counting range, runs the necessary criteria and model constructors, and
locates the first nonzero residual stratum. Here the detailed automated
ledger is rank six. The cited rank-four and rank-five categorification
results are used only in their proved ranges; the larger counting ranges
of the source repository do not by themselves extend those theorems.
